% =============================================================================
% DISCUSSION (compressed for methodological note)
% =============================================================================

\section{Discussion}
\label{sec:discussion}

The value of this paper is methodological. We ask why revenue-based proxies mislead and how widespread the problem is. Our focus is not to estimate the causal effect but to show how proxy substitution alters identification. The \emph{literature audit} (Section~\ref{sec:related}) finds that 83\% (83 of 100) of audited firm-level studies use at least one revenue-related or accounting-based proxy; the issue is prevalent. Journal-tier classification is reported in \ref{app:audit-protocol}. We then show that regressions using revenue rank, ROE, and revenue growth study \emph{revenue dynamics}, not digitalization; same-source bias and construct mismatch make the coefficients uninterpretable for the double-dividend question. The controlled experiment (Section~\ref{sec:synthetic-experiment}, Table~\ref{tab:synthetic-validation}) confirms that significant associations can arise from measurement structure alone when true digitalization is by design uncorrelated with revenue. A reviewer might object that we only demonstrate revenue-based proxy $\to$ revenue-based outcome, which is mechanically correlated. Our point is that many studies in the literature use exactly this structure (revenue rank or similar for digitalization, revenue growth or ROE for outcomes). The controlled experiment shows that even when true digitalization is uncorrelated with revenue, this structure yields significant coefficients; hence results in such designs cannot distinguish a true digitalization effect from same-source structure. We do not claim that all proxy-based designs use this structure, but that the prevalent pattern we document does. We acknowledge that our illustration uses revenue-based outcomes; a demonstration using a non-revenue outcome (e.g., emission or GTFP where data are available) would further strengthen the critique; we encourage such work. Data and code will be deposited in a public repository upon acceptance. We document target vs.\ implemented design, sample flowchart, and merge diagnostics to support a \emph{blueprint} for valid inference: text-based digitalization, LP/OP TFP, emission or GTFP measures, green patents, and regional regulation. Section~\ref{sec:formal-diagnostic} adds a \emph{methodological} component: Proposition~1 formalizes same-source proxy bias; Corollary~1 shows that under monotone same-source proxies the coefficient sign is systematic (rationalizing the negative coefficients in our tables); Theorem~2 shows that the $t$-statistic can diverge with sample size, so false positives become more convincing in larger samples; Corollary~2 states that partial controls (e.g., size or fixed effects) do not eliminate the bias when residual source variation remains. The three-step diagnostic and practical checklist give applied researchers a concrete procedure; our controlled experiment implements the third step.

\textbf{Why significance is not reassurance.} A common empirical intuition is that large samples and robust standard errors provide reassurance. Our formal results show the opposite in same-source settings: if the population projection induced by proxy structure is nonzero (as in Corollary~1), standard asymptotics make the coefficient more, not less, convincingly significant (Theorem~2). In such cases, significance reflects persistence of the source structure rather than evidence on the target causal relationship.

\textbf{Contribution in context.} We provide \emph{formal results} (Proposition~1, Corollary~1, Theorem~2, Corollary~2) and a \emph{usable diagnostic} (three-step procedure plus practical checklist), in addition to the first structured audit of proxy use in this literature, a transparent illustration and controlled experiment, and a blueprint for valid inference. We do not propose a new estimator or identification strategy; the diagnostic allows researchers to assess whether their proxy-based estimates are likely driven by same-source structure. This combination (formal characterization of identification failure, directional bias, and spurious significance, plus diagnostic and demonstration) supports better empirical practice and meets the bar for a method-oriented contribution.

Limitations: the implemented variables do not measure the target constructs; we do not test the double dividend. Total assets match only 64 firm-years (Table~\ref{tab:merge-diagnostics}, Figure~\ref{fig:sample-flow}); we verified the merge and report it for transparency; commercial integrated data (e.g., CSMAR, WIND) could provide larger asset coverage. The IV specification uses a weak instrument (\ref{app:iv}). The sample is Chinese listed firms in heavily polluting industries. The audit covers 100 studies; \ref{app:audit-protocol} reports the protocol and journal-tier classification. The controlled experiment uses a minimal DGP; alternative specifications (e.g., AR(1), industry heterogeneity) could be explored. Some cited works are from 2025--2026 (early view or in press); we cite as of our search date. A valid test would require direct measures and quasi-experimental variation; this note's contribution is to clarify measurement failure and to point to that path.
